Precision medicine leverages patient heterogeneity to estimate individualized treatment regimens—formalized, data-driven approaches designed to match patients with optimal treatments. In the presence of competing events, where multiple causes of failure can occur and one cause precludes others, it is crucial to assess the risk of the specific outcome of interest, such as one type of failure over another. This helps clinicians tailor interventions based on the factors driving that particular cause, leading to more precise treatment strategies. Currently, no precision medicine methods simultaneously account for both survival and competing risk endpoints. To address this gap, we develop a nonparametric individualized treatment regime estimator. Our two-phase method accounts for both overall survival from all events as well as the cumulative incidence of a main event of interest. Additionally, we introduce a multi-utility value function that incorporates both outcomes. We develop random survival and random cumulative incidence forests to construct individual survival and cumulative incidence curves. Simulation studies demonstrated that our proposed method performs well, which we applied to a cohort of peripheral artery disease patients at high risk for limb loss and mortality.